% SGA 3, Expose II, Loop 1 English reconstruction.
% Scope: Proposition 4.5 through Corollary 4.11.8, stopping immediately
% before Section 5.
% Authority: Polo--Gille Exp2-14oct24.pdf, component-local pages 34--43
% (combined-reader pages 92--101; printed pages 82--91).
% OCR was used only as a locator. English wording was drafted independently
% against the French authority and compared with Jacob Reinhold's CC BY 4.0
% English Markdown lineage; that lineage is comparison material, not source
% authority. The PDF controls text, formulas, numbering, notes, and diagrams.

\setcounter{footnote}{70}

% Source page: Exp2 component-local p. 34; combined-reader p. 92; printed p. 82.
\medskip
\noindent\SGAInlineRefTarget{sga3:ii:proposition:4-5}\textbf{命题 4.5.}
\textit{设 \(F\) 是良 \(\mathcal O_S\)-模. 则:}

\textit{(i) \(S\)-群
\(\underline{\operatorname{Aut}}_{\mathcal O_S\text{-mod.}}(F)\)
满足条件 \((E)\), 且有函子性同构}
\[
 \underline{\operatorname{Lie}}\!
   \left(\underline{\operatorname{Aut}}_{\mathcal O_S\text{-mod.}}(F)/S,
         \mathcal M\right)
 \xrightarrow{\sim}
 \underline{\operatorname{Hom}}_{\mathcal O_S\text{-mod.}}
   \left(F,\underline{\operatorname{Lie}}(F/S,\mathcal M)\right),
\]
\textit{它保持由条件 \((E)\) 得到的 \(\mathcal O_S\)-模结构.
特别地, 有 \(\mathcal O_S\)-模同构}
\[
 \underline{\operatorname{Lie}}\!
   \left(\underline{\operatorname{Aut}}_{\mathcal O_S\text{-mod.}}(F)/S\right)
 \xrightarrow{\sim}
 \underline{\operatorname{End}}_{\mathcal O_S\text{-mod.}}(F).
\]

\textit{\SGAInlineRefTarget{sga3:ii:proposition:4-5:item:ii}(ii)\footnote{编者注:
\SGARefLink{sga3:ii:proposition:4-5:item:ii}{(ii)} 项是前面内容的直接推论,
为本版新增.} 此外,
\(\underline{\operatorname{End}}_{\mathcal O_S\text{-mod.}}(F)\)
是良 \(\mathcal O_S\)-模.}

事实上, 由 \SGARefLink{sga3:ii:corollary:4-4-1}{4.4.1}, \(F/S\) 满足 \((E)\), 且
\begin{equation*}
 \underline{\operatorname{Lie}}(F/S,\mathcal M)
 =\underline{\operatorname{Lie}}'(F/S,\mathcal M)
 \simeq
 F\otimes_{\mathcal O_S}
 \underline{\operatorname{Lie}}(\mathbf O_S/S,\mathcal M).
 \tag{1}\label{sga3:ii:proposition:4-5:equation:1}
\end{equation*}
现在由 \SGARefLink{sga3:ii:corollary:4-3-1}{4.3.1} 即得断言 (i). 置
\(E=\underline{\operatorname{End}}_{\mathcal O_S\text{-mod.}}(F)\).
由 (1) 和 \SGARefLink{sga3:ii:remark:4-3-5}{注 4.3.5},
有如下由交换 \(S\)-群态射组成的交换图.

\SGAThreeNativeDiagram[0.90\textwidth]
  {figures/exp2/exp2_p034_endomorphism_differential_square.tex}
  {微分 \(d_E\) 的比较图.}
  {fig:exp2-p45-endomorphism-differential}
  {Polo--Gille Exp2-14oct24.pdf, 分篇第 34 页; 合订本第 92 页; 印刷页码 82}

这里 \(d_F\) 和竖直同构 \((*)\) 都是
\(\mathcal O_S\)-模同构; 因而 \(d_E\) 也是. 这就证明了
\SGARefLink{sga3:ii:proposition:4-5:item:ii}{(ii)}.

\medskip
\noindent\SGAInlineRefTarget{sga3:ii:scholium:4-5-1}\textbf{附论 4.5.1.}
\footnote{编者注: 此附论隐含于原文, 为在
\SGARefLink{sga3:ii:subsection:4-7-1}{4.7.1} 中使用而补出.
这里及下文中的 \(t,t'\) 等表示平方为零的变量;
字母 \(\varepsilon\) 专用于群的单位截面.}
置 (参见 \SGARefLink{sga3:ii:definition:2-1}{2.1})
\(\mathcal O_{I_S}=\mathcal O_S\oplus t\mathcal O_S\), 其中 \(t^2=0\),
并设 \(F\) 是良 \(\mathcal O_S\)-模. 对每个 \(S'\to S\), 态射
\[
 F(S')\oplus tF(S')
 =F(S')\otimes_{\mathcal O(S')}\mathcal O(I_{S'})
 \longrightarrow
 F(I_{S'})
 =F(S')\oplus\underline{\operatorname{Lie}}(F/S)(S'),
\]
在 \(F(S')\) 上为恒等, 因而导出 \(\mathcal O(S')\)-模同构
\[
 tF(S')\xrightarrow{\sim}\underline{\operatorname{Lie}}(F/S)(S').
\]
当 \(S'\) 变动时, 这给出一个同构, 可写为
\[
 \underline{\operatorname{Lie}}(F/S)\simeq tF.
\]

% Source page: Exp2 component-local p. 35; combined-reader p. 93; printed p. 83.
因此, 对每个 \(S'\to S\), \SGARefLink{sga3:ii:proposition:4-5}{命题 4.5}
给出一个交换图.

\SGAThreeNativeDiagram[0.91\textwidth]
  {figures/exp2/exp2_p035_dual_number_automorphism_square.tex}
  {对偶数上的自同态与自同构.}
  {fig:exp2-p451-dual-number-automorphisms}
  {Polo--Gille Exp2-14oct24.pdf, 分篇第 35 页; 合订本第 93 页; 印刷页码 83}

由 4.3 bis 可知, 每个
\[
 X\in
 \underline{\operatorname{End}}_{\mathcal O_{S'}\text{-mod.}}(F_{S'})
\]
对应于
\(\mathrm{id}+tX\in
\underline{\operatorname{Aut}}_{\mathcal O_{I_{S'}}\text{-mod.}}
(F_{I_{S'}})\).

\medskip
\noindent\SGAInlineRefTarget{sga3:ii:definition:4-6}\textbf{定义 4.6.}
取值于群的 \(S\)-函子 \(G\) 称为\emph{良好的 (good)},
若 \(G/S\) 满足条件 \((E)\), 且
\(\underline{\operatorname{Lie}}(G/S)\) 是良
\(\mathcal O_S\)-模.

\footnote{编者注: 下面这句话是新增的.}
注意, 若 \(F\) 是良 \(\mathcal O_S\)-模, 则它是良好
\(S\)-群: 事实上, \(F/S\) 满足 \((E)\), 且
\(\underline{\operatorname{Lie}}(F/S)\simeq F\) 是良
\(\mathcal O_S\)-模.

\medskip
\noindent\SGAInlineRefTarget{sga3:ii:example:4-6-1}\textbf{例 4.6.1.}
\textit{若 \(G\) 可表, 则它是良好的. 事实上, \(G/S\)
满足 \((E)\), 且 \(\underline{\operatorname{Lie}}(G/S)\) 形如
\(V(E)\), 因而由 \SGARefLink{sga3:ii:examples:4-4-2}{4.4.2} 可知它是良模.}

\medskip
\noindent\SGAInlineRefTarget{sga3:ii:lemma:4-6-2}\textbf{引理 4.6.2.}
\footnote{编者注: 此引理是新增的.}
\textit{设 \(G\) 是取值于群的 \(S\)-函子, 且 \(G/S\)
满足 \((E)\), 并置
\(F=\underline{\operatorname{Lie}}(G/S)\). 则 \(F\) 满足
\((E)\), 且交换群态射}
\[
 d:F\longrightarrow\underline{\operatorname{Lie}}(F/S)
\]
\textit{保持 \(\mathcal O_S\)-模结构. 因此, \(G\) 是良好的当且仅当
\(F\to\underline{\operatorname{Lie}}(F/S)\) 是双射.}

\textit{证明.}
假设 \(G/S\) 满足 \((E)\). 设
\(\mathcal M,\mathcal N\) 是有限自由 \(\mathcal O_S\)-模, 记
\[
 F(\mathcal N)=\underline{\operatorname{Lie}}(G/S,\mathcal N),
\]
并用 \(e\) 表示 \(G\) 的单位截面. 对每个 \(S'\to S\),
\[
 F(\mathcal N)(S')
 =
 \left\{
  g\in\operatorname{Hom}_S(I_{S'}(\mathcal N),G)
  \ \middle|\ g\circ\varepsilon_{\mathcal N}=e
 \right\}.
\]
此外,
\(\underline{\operatorname{Lie}}(F(\mathcal N)/S,\mathcal M)(S')\)
与如下态射所成的集合等同:
\[
\left\{
\begin{array}{@{}l@{\ }|@{\ }l@{}}
 \Phi\in
 \operatorname{Hom}_S
   (I_{S'}(\mathcal N)\times_{S'}I_{S'}(\mathcal M),G)
 &
 \begin{aligned}
 \Phi\circ
  (\varepsilon_{\mathcal N}\times
   \mathrm{id}_{I_{S'}(\mathcal M)})
   &=e\in G(I_{S'}(\mathcal M)),\\
 \Phi\circ
  (\mathrm{id}_{I_{S'}(\mathcal N)}\times
   \varepsilon_{\mathcal M})
   &=e\in G(I_{S'}(\mathcal N)).
 \end{aligned}
\end{array}
\right\}.
\]
这就证明了
\begin{equation*}
 \underline{\operatorname{Lie}}(F(\mathcal N)/S,\mathcal M)
 \simeq
 \underline{\operatorname{Lie}}(F(\mathcal M)/S,\mathcal N).
 \tag{1}\label{sga3:ii:lemma:4-6-2:equation:1}
\end{equation*}
由于 \(G/S\) 满足 \((E)\), 可知
\[
\begin{aligned}
 \underline{\operatorname{Lie}}
   (F(\mathcal N)/S,\mathcal M_1\oplus\mathcal M_2)
 &\simeq
 \underline{\operatorname{Lie}}
   (F(\mathcal M_1\oplus\mathcal M_2)/S,\mathcal N)\\
 &\simeq
 \underline{\operatorname{Lie}}
   ((F(\mathcal M_1)\times_SF(\mathcal M_2))/S,\mathcal N).
\end{aligned}
\]
由 \SGARefLink{sga3:ii:proposition:3-8}{3.8}, 最后一项同构于
\[
\begin{aligned}
 &\underline{\operatorname{Lie}}(F(\mathcal M_1)/S,\mathcal N)
 \times_S
 \underline{\operatorname{Lie}}(F(\mathcal M_2)/S,\mathcal N)\\
 &\qquad\simeq
 \underline{\operatorname{Lie}}(F(\mathcal N)/S,\mathcal M_1)
 \times_S
 \underline{\operatorname{Lie}}(F(\mathcal N)/S,\mathcal M_2).
\end{aligned}
\]

% Source page: Exp2 component-local p. 36; combined-reader p. 94; printed p. 84.
因而
\begin{equation*}
 \underline{\operatorname{Lie}}
   (F(\mathcal N)/S,\mathcal M_1\oplus\mathcal M_2)
 \simeq
 \underline{\operatorname{Lie}}(F(\mathcal N)/S,\mathcal M_1)
 \times_S
 \underline{\operatorname{Lie}}(F(\mathcal N)/S,\mathcal M_2),
 \tag{2}\label{sga3:ii:lemma:4-6-2:equation:2}
\end{equation*}
所以由注 4.1.1.2(a) 可知, \(F(\mathcal N)/S\) 满足 \((E)\).

还须证明交换群态射
\[
 d:F(\mathcal N)\longrightarrow
 \underline{\operatorname{Lie}}(F(\mathcal N)/S)
\]
保持 \(\mathcal O_S\)-模结构. 考虑自由
\(\mathcal O_S\)-模 \(\mathcal M=t\mathcal O_S\). 则
\[
 \mathcal O(I_{S'}(\mathcal M))
 =\mathcal O(S')[t]/(t^2)
 =\mathcal O(S')\oplus t\mathcal O(S'),
 \qquad
 \underline{\operatorname{Lie}}(\mathbf O_S/S)\simeq t\mathcal O_S.
\]
设 \(\rho_t:I_S(\mathcal M)\to S\) 是结构态射,
它对应于单射
\[
 u_t:\mathcal O_S\hookrightarrow
 D_{\mathcal O_S}(\mathcal M)
 =\mathcal O_S\oplus t\mathcal O_S.
\]
回忆 (参见 3.4.2), 对每个 \(S\)-函子 \(X\),
\(F(\mathcal N)(X)\) 是满足
\(\phi\circ\varepsilon_{\mathcal N}=e\) 的 \(S\)-态射
\(\phi:I_X(\mathcal N)\to G\) 所成的集合,
而 \(a\in\mathcal O(X)\) 的作用为
\[
 a\cdot\phi=\phi\circ a^*,
\]
其中 \(a^*\) 是与 \(a\) 相应的 \(I_X(\mathcal N)\) 的自同态;
见 \SGARefLink{sga3:ii:subsection:2-1-3}{2.1.3}.

因此, 由 \SGARefLink{sga3:ii:remark:4-3-4}{4.3.4}, 复合
\[
 d:F(\mathcal N)\xrightarrow{\sim}
 F(\mathcal N)\otimes_{\mathcal O_S}t\mathcal O_S
 \longrightarrow
 \underline{\operatorname{Lie}}(F(\mathcal N)/S)
\]
对 \(S'\to S\) 和
\(f\in\operatorname{Hom}_S(S',F(\mathcal N))\) 由下式给出:
\[
 f\longmapsto f\otimes t
 \longmapsto t\cdot(f\circ\rho_t)
 =f\circ\rho_t\circ t^*.
\]
这里 \(f\circ\rho_t\in F(\mathcal N)(I_{S'}(\mathcal M))\)
看作满足 \(\phi\circ\varepsilon_{\mathcal N}=e\) 的 \(S\)-态射
\[
 \phi:I_{I_{S'}(\mathcal M)}(\mathcal N)\longrightarrow G
\]
我们必须证明 \(d\) 是 \(\mathcal O_S\)-线性的, 等价地,
\[
 d(a\cdot f)=u_t(a)\cdot d(f)
 \qquad(a\in\mathcal O(S')).
\]
将 \(f\in F(\mathcal N)(S')\) 看作态射
\(I_{S'}(\mathcal N)\to G\), 同样有
\(a\cdot f=f\circ a^*\). \(\mathcal O(X)\) 在
\(I_X(\mathcal N)\) 上的作用关于 \(X\) 的函子性给出如下交换图.

\SGAThreeNativeDiagram[0.57\textwidth]
  {figures/exp2/exp2_p036_scalar_action_square.tex}
  {标量作用与结构态射 \(\rho_t\) 的相容性.}
  {fig:exp2-p462-scalar-action}
  {Polo--Gille Exp2-14oct24.pdf, 分篇第 36 页; 合订本第 94 页; 印刷页码 84}

所以
\[
\begin{aligned}
 d(a\cdot f)
 &=f\circ a^*\circ\rho_t\circ t^*\\
 &=f\circ\rho_t\circ u_t(a)^*\circ t^*\\
 &=f\circ\rho_t\circ t^*\circ u_t(a)^*
  =u_t(a)\cdot d(f).
\end{aligned}
\]
倒数第二个等号由 \(\mathcal O\) 的交换性得到.
这就完成了 \SGARefLink{sga3:ii:lemma:4-6-2}{引理 4.6.2} 的证明.

\medskip
\noindent\SGAInlineRefTarget{sga3:ii:theorem:4-7}\textbf{定理 4.7.}
\textit{若 \(F\) 是良 \(\mathcal O_S\)-模, 则
\(S\)-群
\(\underline{\operatorname{Aut}}_{\mathcal O_S\text{-mod.}}(F)\)
是良好的.}

\footnote{编者注: 这里利用
\SGARefLink{sga3:ii:corollary:4-3-1}{4.3.1} 和 \SGARefLink{sga3:ii:proposition:4-5}{4.5}
中增加的内容简化了原文.}
事实上, 由 \SGARefLink{sga3:ii:proposition:4-5}{4.5},
\(\underline{\operatorname{Aut}}_{\mathcal O_S\text{-mod.}}(F)/S\)
满足 \((E)\), 而
\[
 \underline{\operatorname{Lie}}\!
  \left(\underline{\operatorname{Aut}}_{\mathcal O_S\text{-mod.}}(F)/S\right)
 \simeq
 \underline{\operatorname{End}}_{\mathcal O_S\text{-mod.}}(F)
\]
是良 \(\mathcal O_S\)-模.

% Source page: Exp2 component-local p. 37; combined-reader p. 95; printed p. 85.
\subsection*{\SGAInlineRefTarget{sga3:ii:subsection:4-7-1}4.7.1.}
\footnote{编者注: 编号 \SGARefLink{sga3:ii:subsection:4-7-1}{4.7.1},
\SGARefLink{sga3:ii:definition:4-7-2}{4.7.2} 和
\SGARefLink{sga3:ii:subsection:4-7-3}{4.7.3} 是新增的.}
设 \(G\) 是 \(S\)-群, \(F\) 是良
\(\mathcal O_S\)-模. 假定给定了 \(G\) 在 \(F\) 中的一个线性表示,
即 (\SGARefLink{sga3:i:definition:4-7-1}{I~4.7.1}) 一个 \(S\)-群态射
\[
 \rho:G\longrightarrow
 \underline{\operatorname{Aut}}_{\mathcal O_S\text{-mod.}}(F).
\]
若 \(G/S\) 满足 \((E)\), 则
\SGARefLink{sga3:ii:remark:4-1-c}{4.1.C} 和 \SGARefLink{sga3:ii:proposition:4-5}{4.5}
给出一个 \(\mathcal O_S\)-模态射, 记为 \(\rho'\) 或 \(d\rho\),
\[
 \underline{\operatorname{Lie}}(G/S)
 \longrightarrow
 \underline{\operatorname{Lie}}\!
  \left(\underline{\operatorname{Aut}}_{\mathcal O_S\text{-mod.}}(F)/S\right)
 \simeq
 \underline{\operatorname{End}}_{\mathcal O_S\text{-mod.}}(F).
\]
\footnote{编者注: 下面这句话是新增的.}
此外, 置
\(\mathcal O_{I_S}=\mathcal O_S\oplus t\mathcal O_S\), 其中
\(t^2=0\). 由 \SGARefLink{sga3:ii:scholium:4-5-1}{4.5.1} 可知, 若 \(S'\to S\) 且
\[
 X\in\underline{\operatorname{Lie}}(G/S)(S')\subset G(I_{S'}),
\]
则在
\(\underline{\operatorname{Aut}}_{\mathcal O_{I_{S'}}\text{-mod.}}
(F_{I_{S'}})\) 中有
\begin{equation*}
 \rho(X)=\mathrm{id}+t\rho'(X).
 \tag{*}\label{sga3:ii:subsection:4-7-1:equation:star}
\end{equation*}
换言之, 对每个 \(S''\to I_{S'}\) 和 \(f\in F(S'')\),
\[
 \rho(X)(f)=f+t\rho'(X)(f)
 \quad\text{于 }F(S'').
\]

\medskip
\noindent\SGAInlineRefTarget{sga3:ii:definition:4-7-2}\textbf{定义 4.7.2.}
设 \(G\) 是良好 \(S\)-群. 则
\(\underline{\operatorname{Lie}}(G/S)\) 是良
\(\mathcal O_S\)-模, 且有 \(S\)-群态射
\[
 \operatorname{Ad}:G\longrightarrow
 \underline{\operatorname{Aut}}_{\mathcal O_S\text{-mod.}}
   \bigl(\underline{\operatorname{Lie}}(G/S)\bigr).
\]
由 \SGARefLink{sga3:ii:subsection:4-7-1}{4.7.1}, 它导出一个 \(\mathcal O_S\)-模态射
\[
 \operatorname{ad}:
 \underline{\operatorname{Lie}}(G/S)
 \longrightarrow
 \underline{\operatorname{End}}_{\mathcal O_S\text{-mod.}}
   \bigl(\underline{\operatorname{Lie}}(G/S)\bigr),
\]
等价地, 一个 \(\mathcal O_S\)-双线性态射
\[
 \underline{\operatorname{Lie}}(G/S)
 \times_S
 \underline{\operatorname{Lie}}(G/S)
 \longrightarrow
 \underline{\operatorname{Lie}}(G/S),
 \qquad
 (x,y)\longmapsto[x,y]=\operatorname{ad}(x)\cdot y.
\]
这里 \(x,y\) 是
\[
 \underline{\operatorname{Lie}}(G/S)(S')
 =
 \underline{\operatorname{Lie}}(G_{S'}/S')(S').
\]
中的任意元素. 若 \(G\) 交换, 则 \([x,y]=0\).

\subsection*{\SGAInlineRefTarget{sga3:ii:subsection:4-7-3}4.7.3.}
括号运算还有如下等价描述. 只需考虑 \(S\) 上的截面, 即
\[
 x,y\in\underline{\operatorname{Lie}}(G/S)(S).
\]
有典范同构
\(I_S\times_SI_S\simeq I_{I_S}\). 为避免混淆, 用 \(I\) 和
\(I'\) 表示 \(I_S\) 的两个副本, 并置
\[
 \mathcal O_I=\mathcal O_S[t],
 \qquad
 \mathcal O_{I'}=\mathcal O_S[t'],
 \qquad
 t^2=0=(t')^2.
\]
于是有交换方块

\SGAThreeNativeDiagram[0.32\textwidth]
  {figures/exp2/exp2_p037_dual_number_base_square.tex}
  {\(I\times_S I'\) 上的两个对偶数结构.}
  {fig:exp2-p473-dual-number-base}
  {Polo--Gille Exp2-14oct24.pdf, 分篇第 37 页; 合订本第 95 页; 印刷页码 85}

其中从 \(I\times I'\) 出发的两个箭头分别将其与
\(I\) 上或 \(I'\) 上的对偶数概形等同.
记 \(L=\underline{\operatorname{Lie}}(G/S)\), 得到如下群的交换图.

% Source page: Exp2 component-local p. 38; combined-reader p. 96; printed p. 86.
\SGAThreeNativeDiagram[0.69\textwidth]
  {figures/exp2/exp2_p038_exact_sequence_diagram_1.tex}
  {用于构造括号运算的第一个正合列图.}
  {fig:exp2-p473-exact-sequences-one}
  {Polo--Gille Exp2-14oct24.pdf, 分篇第 38 页; 合订本第 96 页; 印刷页码 86}

此图中的第九个对象是
\(\underline{\operatorname{Lie}}(L/S)(S)\). 若 \(G\) 是良好的,
它就是 \(L(S)\), 因而得到下一个交换图. 其行和列都是群的正合列;
记为 \(L(-)\) 的五个群都是交换的.
\footnote{编者注: 下面扩充了原文.} 在等同
\[
 L(I)=L(S)\oplus tL(S),
 \qquad
 L(I')=L(S)\oplus t'L(S),
\]
之下, 单射 \(L(S)\hookrightarrow L(I)\) 和
\(L(S)\hookrightarrow L(I')\) 分别为 \(u\mapsto tu\) 和
\(u\mapsto t'u\).

\SGAThreeNativeDiagram[0.69\textwidth]
  {figures/exp2/exp2_p038_exact_sequence_diagram_2.tex}
  {定义 \([x,y]\) 的正合列图.}
  {fig:exp2-p473-exact-sequences-two}
  {Polo--Gille Exp2-14oct24.pdf, 分篇第 38 页; 合订本第 96 页; 印刷页码 86}

在这样的图中, 如图取 \(x,y\), 并任取
\(x\) 的提升 \(\widetilde x\in L(I)\) 和 \(y\) 的提升 \(y'\in L(I')\).
换位子
\[
 \widetilde x\,y'\,\widetilde x^{-1}(y')^{-1}
 \quad\text{于 }G(I\times I')
\]
与所选提升无关, 并且如图所示, 它是一个唯一确定的元素 \(z\) 的像.
有 \(z=[x,y]\).
\footnote{编者注: 原文的一条脚注说, 作者应 Gabriel 的要求承认,
这个习题并非显然, 而这也正是没有将它放入正文的原因.
这里考虑了 \SGARefLink{sga3:ii:subsection:4-7-1}{4.7.1} 中的增补,
给出了细节; 另见 [DG70], \(\S\)II.4, \SGARefLink{sga3:ii:subsection:4-2}{4.2}.}

事实上, 对 \(x\) 和 \(y\) 在典范截面
\(L(S)\to L(I)\) 和 \(L(S)\to L(I')\) 下的像仍用相同的字母表示.
则
\[
 \widetilde x=xu,\qquad y'=yv,
 \qquad u,v\in L(S)=L(I)\cap L(I').
\]
由于 \(L(I)\) 和 \(L(I')\) 交换,
\[
\begin{aligned}
 \widetilde x\,y'\,\widetilde x^{-1}(y')^{-1}
 &=xuyvu^{-1}x^{-1}v^{-1}y^{-1}\\
 &=xuyu^{-1}vx^{-1}v^{-1}y^{-1}\\
 &=xyx^{-1}y^{-1}.
\end{aligned}
\]

% Source page: Exp2 component-local p. 39; combined-reader p. 97; printed p. 87.
此元素在 \(G(I)\) 和 \(G(I')\) 中的像都是单位元,
因此来自上面所示的唯一的 \(z\). 最后, 将 \(y\) 看成
\(L_{I'}(I')\) 的元素, 并将 \(x\) 看成
\(L(S)\subset G(I')\) 的元素, 则 \SGARefLink{sga3:ii:subsection:4-7-1}{公式 4.7.1}
\((*)\) 给出
\[
 xyx^{-1}
 =\operatorname{Ad}(x)(y)
 =(\mathrm{id}+t'\operatorname{ad}(x))(y)
 =y+t'[x,y].
\]
所以 \(xyx^{-1}y^{-1}\in L(I')\) 是
\(z=[x,y]\in L(S)\) 的像.

从这一构造现在可清楚看出两个性质:

\textit{\SGAInlineRefTarget{sga3:ii:subsection:4-7-3:item:i}(i)}
括号运算关于 \(G\) 是函子性的. 更确切地说,
\[
 G\longmapsto\underline{\operatorname{Lie}}(G/S)
\]
是从良好 \(S\)-群到带有一个 \(\mathcal O_S\)-双线性合成律的
良 \(\mathcal O_S\)-模的函子;

\textit{(ii)} \([x,y]+[y,x]=0\), 因为此图关于其主对角线对称.
\footnote{编者注: 即 \(x\) 和 \(y\) 的地位对称.
\([y,x]\) 在 \(G(I\times I')\) 中的像是换位子
\(y'\widetilde x(y')^{-1}\widetilde x^{-1}\), 它是
\(\widetilde x y'\widetilde x^{-1}(y')^{-1}=[x,y]\) 的逆元.
另一方面 (参见下面的 \SGARefLink{sga3:ii:definition:4-10}{4.10}),
此时尚不知道是否必有 \([x,x]=0\): 若 \(x\) 由
\(\widetilde x\in L(I)\) 及 \(x'\in L(I')\) 表示,
则先验地并不清楚 \(\widetilde x\) 和 \(x'\) 是否交换.}

\medskip
\noindent\SGAInlineRefTarget{sga3:ii:proposition:4-8}\textbf{命题 4.8.}
\textit{设 \(F\) 是良 \(\mathcal O_S\)-模. 在等同}
\[
 \underline{\operatorname{Lie}}\!
  \left(\underline{\operatorname{Aut}}_{\mathcal O_S\text{-mod.}}(F)/S\right)
 =
 \underline{\operatorname{End}}_{\mathcal O_S\text{-mod.}}(F),
\]
\textit{之下, 有}
\[
 \operatorname{Ad}(g)\cdot Y=g\circ Y\circ g^{-1},
 \qquad
 [X,Y]=X\circ Y-Y\circ X,
\]
\textit{其中 \(S'\to S\),
\(g\in
\underline{\operatorname{Aut}}_{\mathcal O_{S'}\text{-mod.}}(F_{S'})\)
以及}
\[
 X,Y\in
 \underline{\operatorname{Lie}}\!
  \left(\underline{\operatorname{Aut}}_{\mathcal O_S\text{-mod.}}(F)/S\right)
  (S')
 =
 \underline{\operatorname{End}}_{\mathcal O_{S'}\text{-mod.}}(F_{S'}).
\]
\textit{都是任意的.}

\textit{证明.}\footnote{编者注: 下面扩充并简化了原文.}
作基变换后, 可假设 \(S'=S\). 置 \(I=I_S\) 及
\(\mathcal O_I=\mathcal O_S[t]\), 其中 \(t^2=0\).
回忆 \SGARefLink{sga3:ii:scholium:4-5-1}{4.5.1} 中的包含映射
\[
 i:\underline{\operatorname{End}}_{\mathcal O_S\text{-mod.}}(F)
 \hookrightarrow
 \underline{\operatorname{Aut}}_{\mathcal O_I\text{-mod.}}(F_I)
\]
将 \(Y\) 映到 \(\mathrm{id}+tY\). 由 4.1.A 中
\(\operatorname{Ad}(g)\) 的定义,
\[
 \mathrm{id}+t\operatorname{Ad}(g)(Y)
 =g\circ(\mathrm{id}+tY)\circ g^{-1}
 =\mathrm{id}+t(g\circ Y\circ g^{-1}),
\]
因而
\(\operatorname{Ad}(g)(Y)=g\circ Y\circ g^{-1}\).

设 \(I'\) 是 \(I_S\) 的第二个副本, 并置
\(\mathcal O_{I'}=\mathcal O_S[t']\), 其中 \((t')^2=0\).
将 \SGARefLink{sga3:ii:subsection:4-7-3}{4.7.3} 应用于
\[
 G=\underline{\operatorname{Aut}}_{\mathcal O_S\text{-mod.}}(F),
 \qquad
 L=\underline{\operatorname{Lie}}(G/S)
  =\underline{\operatorname{End}}_{\mathcal O_S\text{-mod.}}(F).
\]
这里 \(L=\operatorname{End}(F)\) 是
\SGARefLink{sga3:ii:proposition:4-5}{命题 4.5} 中的微分等同.
重版本在此处印作 \(\operatorname{Aut}(F)\), 但
\(\operatorname{End}(F)\) 才是下文自同态计算所用的类型正确的对象.
将 \(X\) 与其在典范截面 \(L(S)\hookrightarrow L(I)\) 下的像等同.
它在 \(G(I\times I')\) 中的像是 \(\mathrm{id}+t'X\),
其逆元为 \(\mathrm{id}-t'X\). 同样, \(Y\) 提升为
\(\mathrm{id}+tY\), 其逆元为 \(\mathrm{id}-tY\). 它们的换位子为
\[
\begin{aligned}
 &(\mathrm{id}+t'X)\circ(\mathrm{id}+tY)
 \circ(\mathrm{id}-t'X)\circ(\mathrm{id}-tY)\\
 &\hspace{5em}
 =\mathrm{id}+tt'(X\circ Y-Y\circ X).
\end{aligned}
\]
这是 \(Z=X\circ Y-Y\circ X\in L(S)\) 在 \(G(I\times I')\)
中的像: 首先 \(Z\) 映到 \(tZ\in L(I)\), 再映到
\(\mathrm{id}+t'tZ\in G(I\times I')\).
由 \SGARefLink{sga3:ii:subsection:4-7-3}{4.7.3},
\([X,Y]=X\circ Y-Y\circ X\).

% Source page: Exp2 component-local p. 40; combined-reader p. 98; printed p. 88.
\medskip
\noindent\SGAInlineRefTarget{sga3:ii:corollary:4-8-1}\textbf{推论 4.8.1.}
\textit{设 \(G\) 是良好 \(S\)-群, 且
\(x,y,z\in\underline{\operatorname{Lie}}(G/S)(S')\). 则}
\[
 [x,[y,z]]+[y,[z,x]]+[z,[x,y]]=0.
\]
\footnote{编者注: 下面扩充了原文.}
事实上, 由于 \(G\) 是良好的,
\(\underline{\operatorname{Lie}}(G/S)\) 是良
\(\mathcal O_S\)-模; 因此 \SGARefLink{sga3:ii:theorem:4-7}{定理 4.7} 表明
\[
 \underline{\operatorname{Aut}}_{\mathcal O_S\text{-mod.}}
   \bigl(\underline{\operatorname{Lie}}(G/S)\bigr)
\]
是良好 \(S\)-群. 态射
\[
 \operatorname{Ad}:G\longrightarrow
 \underline{\operatorname{Aut}}_{\mathcal O_S\text{-mod.}}
  \bigl(\underline{\operatorname{Lie}}(G/S)\bigr)
\]
与 \SGARefLink{sga3:ii:subsection:4-7-3:item:i}{4.7.3(i)} 中的函子性给出
\[
 \operatorname{ad}[x,y]
 =[\operatorname{ad}x,\operatorname{ad}y].
\]
结合 \SGARefLink{sga3:ii:proposition:4-8}{命题 4.8}, 得到
\[
 \operatorname{ad}[x,y]
 =\operatorname{ad}x\circ\operatorname{ad}y
  -\operatorname{ad}y\circ\operatorname{ad}x.
\]
将两端作用于 \(z\), 就得到 Jacobi 恒等式.

\medskip
\noindent\SGAInlineRefTarget{sga3:ii:corollary:4-8-2}\textbf{推论 4.8.2.}
\textit{设良好 \(S\)-群 \(G\) 线性作用于良
\(\mathcal O_S\)-模 \(F\); 等价地, 设 \(F\) 是
\(G\)-\(\mathcal O_S\)-模, 且 \(G\) 和 \(F\) 分别是良好群和良模.
则线性映射}
\[
 \rho':\underline{\operatorname{Lie}}(G/S)
 \longrightarrow
 \underline{\operatorname{End}}_{\mathcal O_S\text{-mod.}}(F)
\]
\textit{是一个表示; 即}
\[
 \rho'([x,y])
 =\rho'(x)\circ\rho'(y)-\rho'(y)\circ\rho'(x).
\]

\medskip
\noindent\SGAInlineRefTarget{sga3:ii:scholium:4-9}\textbf{附论 4.9.}
对每个良好 \(S\)-群 \(G\) (例如每个可表的群),
我们已函子性地配给一个良 \(\mathcal O_S\)-模
\(\underline{\operatorname{Lie}}(G/S)\), 它带有满足
\[
 [x,y]+[y,x]=0,
 \qquad
 [x,[y,z]]+[y,[z,x]]+[z,[x,y]]=0.
\]
的双线性映射. 我们将带有这一结构的
\(\underline{\operatorname{Lie}}(G/S)\) 称为 \(G\) 在 \(S\) 上的
\emph{Lie 代数}. 这里需要加以限定: 尚不知道它是否是严格意义下的
\(\mathcal O_S\)-Lie 代数.
\footnote{编者注: 未解决的一点是是否有 \([x,x]=0\);
见 \SGARefLink{sga3:ii:definition:4-10}{4.10}.}
\(G\) 在良 \(\mathcal O_S\)-模 \(F\) 中的每个线性表示都导出其
Lie 代数的一个表示. 特别地, \(G\) 的伴随表示导出其 Lie 代数的伴随表示.

\medskip
\noindent\SGAInlineRefTarget{sga3:ii:definition:4-10}\textbf{定义 4.10.}
取值于群的 \(S\)-函子 \(G\) 称为\emph{甚好的 (very good)},
若它是良好的, 且 \(\underline{\operatorname{Lie}}(G/S)\) 是
\(\mathcal O_S\)-Lie 代数, 即恒有 \([x,x]=0\).

\medskip
\noindent\SGAInlineRefTarget{sga3:ii:example:4-10-1}\textbf{例 4.10.1.}
下列 \(S\)-群是甚好的:
每个良 \(\mathcal O_S\)-模 \(F\) 的
\(\underline{\operatorname{Aut}}_{\mathcal O_S\text{-mod.}}(F)\)
(由 \SGARefLink{sga3:ii:theorem:4-7}{4.7} 和 \SGARefLink{sga3:ii:proposition:4-8}{4.8});
每个可表的群 (将在下面证明); 每个容许到某个甚好 \(S\)-群的单态射的
良好 \(S\)-群, 特别是可表群的每个良好子群函子; 以及每个在良
\(\mathcal O_S\)-模中容许忠实线性表示的良好 \(S\)-群,
例如每个 \(\operatorname{Ad}\) 为单态射的良好 \(S\)-群.

% Source page: Exp2 component-local p. 41; combined-reader p. 99; printed p. 89.
\subsection*{\SGAInlineRefTarget{sga3:ii:subsection:4-11}4.11.}
现在假设 \(G\) 是 \(S\) 上的群概形. 由 4.1.4,
\(\underline{\operatorname{Lie}}(G/S)(S)\) 与 \(G/S\)
的在右平移下不变的无穷小自同构所成的群等同;
由 \SGARefLink{sga3:ii:subsection:3-14}{3.14}, 它就是
\(\mathcal O_G\) 在 \(\mathcal O_S\) 上的在右平移下不变的导子所成的群.
此等同保持模结构, 并且是 Lie 代数的反同构.
\footnote{编者注: 这里改正了原文的一个符号错误;
参见 [DG70], \(\S\)II.4, 4.4 和 \SGARefLink{sga3:ii:definition:4-6}{4.6}.}

为说明这一点, 如 \SGARefLink{sga3:ii:subsection:4-7-3}{4.7.3} 中那样论证. 置
\[
 \mathcal O_I=\mathcal O_S[t],
 \qquad
 \mathcal O_{I'}=\mathcal O_S[t'],
 \qquad
 t^2=0=(t')^2,
\]
并取 \(x\in L(I)\), \(y\in L(I')\), 其中
\(L=\underline{\operatorname{Lie}}(G/S)\). 左平移
\(\lambda_x\), 分别地, \(\lambda_y\), 是
\(G_{I\times I'}\) 的 \(S\)-自同构, 它在 \(G_{I'}\),
分别地, 在 \(G_I\) 上诱导恒等. 它对应于
\[
 \mathcal O_{G_{I\times I'}}
 =\mathcal O_G[t,t']/(t^2,(t')^2),
\]
的 \(\mathcal O_S\)-自同构
\[
 u=\mathrm{id}+t\,d_x,
 \qquad\text{分别地}\qquad
 v=\mathrm{id}+t'\,d_y
\]
其中 \(d_x,d_y\) 是 \(\mathcal O_G\) 的在右平移下不变的
\(\mathcal O_S\)-导子. \(G_{I\times I'}\) 的 \(S\)-自同构与
\(\mathcal O_{G_{I\times I'}}\) 的 \(\mathcal O_S\)-自同构之间的对应
是反变的. 因此
\[
 \lambda_x\lambda_y\lambda_x^{-1}\lambda_y^{-1}
\]
对应于
\[
 v^{-1}u^{-1}vu
 =\mathrm{id}+tt'(d_yd_x-d_xd_y).
\]
由 \SGARefLink{sga3:ii:subsection:4-7-3}{4.7.3} 可知,
\[
 x\longmapsto-d_x
\]
是 Lie 代数同构; 更多细节见 [DG70],
\(\S\)II.4, 4.4 和 \SGARefLink{sga3:ii:definition:4-6}{4.6}.
对每个 \(S'\to S\), 同一论证也适用于
\[
 \underline{\operatorname{Lie}}(G/S)(S')
 =
 \underline{\operatorname{Lie}}(G_{S'}/S')(S')
\]
于是重新得到了经典定义.
\footnote{编者注: 下面扩充了原文. 特别是为便于后文引用,
新增了编号 \SGARefLink{sga3:ii:scholium:4-11-1}{4.11.1}--\SGARefLink{sga3:ii:corollary:4-11-8}{4.11.8}.}

\medskip
\noindent\SGAInlineRefTarget{sga3:ii:scholium:4-11-1}\textbf{附论 4.11.1.}
\textit{通过同构 \(x\mapsto-d_x\),
\(\underline{\operatorname{Lie}}(G/S)\) 与如下函子等同:
它将每个 \(S'\to S\) 映到 \(G_{S'}\) 相对于 \(S'\)
的在右平移下不变的导子所成的 \(\mathcal O(S')\)-Lie 代数.}

由 \SGARefLink{sga3:ii:example:4-6-1}{4.6.1}, 已知每个可表的群是良好的,
因此前面的讨论给出:

\medskip
\noindent\SGAInlineRefTarget{sga3:ii:corollary:4-11-2}\textbf{推论 4.11.2.}
\textit{每个可表的群都是甚好的.}

设 \(\varepsilon:S\to G\) 是 \(G\) 的单位截面, 并置
\[
 \omega^1_{G/S}=\varepsilon^*(\Omega^1_{G/S}).
\]
回忆 \SGARefLink{sga3:ii:proposition:3-3}{3.3},
\(\underline{\operatorname{Lie}}(G/S)\) 由向量纤维空间
\(V(\omega^1_{G/S})\) 表示.

\medskip
\noindent\SGAInlineRefTarget{sga3:ii:scholium:4-11-3}\textbf{附论 4.11.3.}
\textit{因此, 对每个 \(S\)-群概形 \(G\), 我们已函子性地配给一个
\(S\) 上的向量纤维空间}
\[
 \operatorname{Lie}(G/S)=V(\omega^1_{G/S})
\]
\textit{它表示函子 \(\underline{\operatorname{Lie}}(G/S)\),
从而带有取值于 \(\mathcal O_S\)-Lie 代数的 \(S\)-概形结构.
此外, 此构造与基扩张及有限积可交换
(参见 \SGARefLink{sga3:ii:proposition:3-4}{3.4} 和 \SGARefLink{sga3:ii:proposition:3-8}{3.8}).}

\medskip
\noindent\SGAInlineRefTarget{sga3:ii:remark:4-11-4}\textbf{注 4.11.4.}
\footnote{编者注: 下面的内容考虑了第 I 篇
\(\S\)\SGARefLink{sga3:i:subsection:6-8}{6.8} 中的增补而作了扩充.}
设 \(\pi:G\to S\) 是结构态射.

\textit{(a)} \(\mathcal O_G\)-模 \(\Omega^1_{G/S}\) 自然地是
\((G\times_SG)\)-等变的 (参见 I, \(\S\)6). 因而由
\SGARefLink{sga3:i:proposition:6-8-1}{I~6.8.1},
\[
 \Omega^1_{G/S}\simeq\pi^*(\omega^1_{G/S}).
\]
例如, 若 \(\omega^1_{G/S}\) 局部自由, 分别地, 局部自由且为有限型,
则 \(\Omega^1_{G/S}\) 也如此. 特别地, 当 \(S\) 是域的谱,
分别地, 当 \(S\) 是域的谱且 \(G\) 在 \(S\) 上局部有限型时,
就有这种情形.

% Source page: Exp2 component-local p. 42; combined-reader p. 100; printed p. 90.
\textit{(b)} 由 \SGARefLink{sga3:i:remark:6-8-2}{I~6.8.2},
\(\omega^1_{G/S}\) 有典范的
\(G\)-\(\mathcal O_S\)-模结构, 它导出
\[
 V(\omega^1_{G/S})=\operatorname{Lie}(G/S).
\]
上的伴随作用.
\footnote{编者注: 因此, \(G\) 在 \(\omega^1_{G/S}\) 上的线性作用
可称为前伴随作用; 由于术语的滥用, 我们仍称它为伴随作用.
第 III 篇 (特别是 \SGARefLink{sga3:iii:paragraph:0-8}{III~0.8})
中使用的一个稍一般的构造如下. 假定对每个 \(Y\to S\), 配给一个
\(\mathcal O_Y\)-模 \(\mathcal N(Y)\),
并对每个 \(S\)-态射 \(\phi:Z\to Y\), 函子性地配给一个
\(\mathcal O_Z\)-模态射
\(\mathcal N(\phi):\phi^*\mathcal N(Y)\to\mathcal N(Z)\).
例如, 若 \(\mathcal J\) 是 \(\mathcal O_S\) 的一个拟凝聚理想,
则可取 \(\mathcal N(Y)=\mathcal J\mathcal O_Y\).
\(G\) 在 \(\omega^1_{G/S}\) 上的作用便导出 \(G\)
在一个取值于交换群的 \(S\)-函子上的广义伴随作用,
该函子将 \(f:Y\to S\) 映到群
\(\operatorname{Hom}_{\mathcal O_Y}
(f^*\omega^1_{G/S},\mathcal N(Y))\).}

\textit{(c)} 由 EGA~I, 5.3.11 和 5.4.6,
\(\varepsilon:S\to G\) 是浸入; 若 \(G\) 在 \(S\) 上分离,
则它是闭浸入. 因此, \(\omega^1_{G/S}\) 与
\(\mathcal I/\mathcal I^2\) 等同, 其中 \(\mathcal I\)
是在某个开子概形 \(U\subset G\) 中定义 \(\varepsilon(S)\) 的拟凝聚理想,
并且 \(\varepsilon(S)\) 在其中闭.
若 \(G\to S\) 分离, 则可取 \(U=G\).
若 \(G=\operatorname{Spec}A(G)\) 在 \(S\) 上仿射,
则 \(\mathcal I\) 是 \(A(G)\) 的增广理想, 即
\[
 \varepsilon^\natural:A(G)\longrightarrow\mathcal O_S
\]
的核 (参见 \SGARefLink{sga3:i:subsection:4-2}{I~4.2}).
\footnote{编者注: 前面关于 \(\omega^1_{G/S}\) 的描述是新增的,
将在后文用到 (例如 VII\textsubscript{A}, \SGARefLink{sga3:VIIA:5.5}{5.5}).}

\medskip
\noindent\SGAInlineRefTarget{sga3:ii:remark:4-11-5}\textbf{注 4.11.5.}
同构
\(\Omega^1_{G/S}\simeq\pi^*(\omega^1_{G/S})\) 将
\(\mathcal O_S\)-模 \(\omega^1_{G/S}\) 与 \(G\) 相对于 \(S\)
的在右平移下不变的微分所成的层
\[
 \pi_*(\Omega^1_{G/S})^G
\]
等同: 它在开集 \(U\subset S\) 上的截面就是
\(\Omega^1_{G/S}\) 在 \(\pi^{-1}(U)\) 上的在右平移下不变的截面
(参见 \SGARefLink{sga3:i:remark:6-8-3}{I~6.8.3}; 亦可比较
VII\textsubscript{A}, \SGARefLink{sga3:VIIA:2.4}{2.4}).
\footnote{编者注: 下面不使用这一借助不变微分对
\(\omega^1_{G/S}\) 的描述.}

\medskip
\noindent\SGAInlineRefTarget{sga3:ii:notation:4-11-6}\textbf{记号 4.11.6.}
用 \(\mathop{\mathscr L\!ie}(G/S)\) 表示向量纤维空间
\(\operatorname{Lie}(G/S)\to S\) 的截面层. 它是
\(\mathcal O_S\)-模
\[
 \mathop{\mathscr L\!ie}(G/S)
 =(\omega^1_{G/S})^\vee
 =\mathcal Hom_{\mathcal O_S}(\omega^1_{G/S},\mathcal O_S),
\]
与 \(\omega^1_{G/S}\) 对偶 (参见 EGA~II, 1.7.9),
并且带有一个 \(\mathcal O_S\)-Lie 代数结构.

一般而言, 此构造与基扩张不可交换. 因此, 此模上的 Lie 代数结构并不确定
\(\operatorname{Lie}(G/S)\) 上的取值于 \(\mathcal O_S\)-Lie 代数的
\(S\)-概形结构.
\footnote{编者注: 原因是
\(\mathop{\mathscr L\!ie}(G/S)\simeq(\omega^1_{G/S})^\vee\)
未必确定 \(\omega^1_{G/S}\). 例如, 设
\(S=\mathbb A^1_k\), 其中 \(k\) 是域,
并设 \(G\) 是在 \(s=0\) 处的纤维为 \(\mathbb G_{a,k}\),
而其他所有纤维均为单位群的 \(S\)-群. 则
\(\mathop{\mathscr L\!ie}(G/S)=0\), 但
\(\operatorname{Lie}(G_s/k)(k)=k\).}

不过:

\medskip
\noindent\SGAInlineRefTarget{sga3:ii:lemma:4-11-7}\textbf{引理 4.11.7.}
\textit{假设 \(\omega^1_{G/S}\) 局部自由且为有限型.
特别地, 若 \(G\) 在 \(S\) 上光滑
(参见 EGA~IV\textsubscript{4}, 17.2.4), 或者若 \(S\) 是域的谱
且 \(G\) 在 \(S\) 上局部有限型, 就有这种情形. 则}
\[
 \mathop{\mathscr L\!ie}(G/S)^\vee
 \simeq(\omega^1_{G/S})^{\vee\vee}
 \simeq\omega^1_{G/S},
\]
\textit{从而}
\[
 \operatorname{Lie}(G/S)
 =V(\omega^1_{G/S})
 =V\bigl(\mathop{\mathscr L\!ie}(G/S)^\vee\bigr)
 =W\bigl(\mathop{\mathscr L\!ie}(G/S)\bigr),
\]
\textit{其中最后一个等号由 \SGARefLink{sga3:i:corollary:4-6-5-1}{I~4.6.5.1} 得到.}

% Source page: Exp2 component-local p. 43; combined-reader p. 101; printed p. 91.
最后, 设 \(G\to H\) 是取值于群的函子的单态射. 则
\[
 \operatorname{Lie}(G/S)\longrightarrow\operatorname{Lie}(H/S)
\]
也是单态射 (参见 \SGARefLink{sga3:ii:proposition:3-7}{3.7}).
向量纤维空间之间的每个向量单态射都是闭浸入.
\footnote{编者注: 设
\(f:M\to N\) 是 \(\mathcal O_S\)-模态射,
并设 \(P=\operatorname{Coker}(f)\). 若 \(V(N)\to V(M)\) 是单态射,
则满态射
\(\operatorname{Sym}(N)\to\operatorname{Sym}(P)\) 经过
\(\mathcal O_S\) 分解; 因而 \(P=0\).}
于是得到:

\medskip
\noindent\SGAInlineRefTarget{sga3:ii:corollary:4-11-8}\textbf{推论 4.11.8.}
\textit{设 \(G\hookrightarrow H\) 是 \(S\)-群概形的单态射.}

\textit{(i) 态射}
\[
 \operatorname{Lie}(G/S)\longrightarrow\operatorname{Lie}(H/S)
\]
\textit{是闭浸入; 因此}
\[
 \omega^1_{H/S}\longrightarrow\omega^1_{G/S}
\]
\textit{是满态射.}

\textit{(ii) 若 \(\omega^1_{G/S}\) 局部自由且为有限型,
则相应的态射}
\[
 \mathop{\mathscr L\!ie}(G/S)\longrightarrow
 \mathop{\mathscr L\!ie}(H/S)
\]
\textit{将 \(\mathop{\mathscr L\!ie}(G/S)\) 与
\(\mathop{\mathscr L\!ie}(H/S)\) 的一个局部为直和项的
\(\mathcal O_S\)-子模等同.}

% Section 5 begins immediately after this point on source local p. 43
% and is intentionally excluded from this component.

% Translation coverage: Proposition 4.5 through Corollary 4.11.8, before Section 5.
% Source lines: 972; translated body lines before this audit comment: 952.
% Math preserved: 411 inline formulas, 73 display blocks, 4 equation* environments; prose \text arguments translated.
% Footnotes: 21/21; reference targets: 26/26; reference links: 59/59; equation labels: 4/4; diagrams: 6/6.
% Unresolved TODOs: 0; compile and visual review delegated to the volume integrator.
